% !TEX program = lualatex
\documentclass[lualatex,ja=standard]{bxjsarticle}

\usepackage{xcolor}   % テキストカラー
\usepackage{ulem}     % 取り消し線・波線下線

\title{テキスト装飾のサンプル}
\author{LaTeX サンプル}
\date{\today}

\begin{document}

\maketitle

% -------------------------------------------------------------------
% 基本的な文字装飾
% -------------------------------------------------------------------
\section{基本装飾}

\textbf{太字 (bold)}

\textit{イタリック (italic)}

\underline{下線 (underline)}

\emph{強調 (emph)} % 通常文中はイタリック、イタリック文中は直立体

\sout{取り消し線 (strikethrough)} % ulem パッケージ必要

\uwave{波線下線 (wavy underline)} % ulem パッケージ必要

% -------------------------------------------------------------------
% フォントファミリー
% -------------------------------------------------------------------
\section{フォントファミリー}

\textrm{ローマン体 (roman)} % デフォルト

\textsf{サンセリフ体 (sans-serif)}

\texttt{等幅フォント (typewriter / monospace)}

% -------------------------------------------------------------------
% フォントサイズ
% -------------------------------------------------------------------
\section{フォントサイズ}

{\tiny       tiny}
{\scriptsize scriptsize}
{\footnotesize footnotesize}
{\small      small}
{\normalsize normalsize（標準）}
{\large      large}
{\Large      Large}
{\LARGE      LARGE}
{\huge       huge}
{\Huge       Huge}

% -------------------------------------------------------------------
% テキストカラー（xcolor パッケージ）
% -------------------------------------------------------------------
\section{テキストカラー}

\textcolor{red}{赤}
\textcolor{blue}{青}
\textcolor{green}{緑}
\textcolor{orange}{オレンジ}
\textcolor{purple}{紫}

% RGB で任意の色を指定
\textcolor[RGB]{0,128,128}{カスタムカラー (teal)}

% 背景色
\colorbox{yellow}{背景色: 黄色}
\colorbox{cyan}{\textcolor{white}{背景色: シアン + 白文字}}

% -------------------------------------------------------------------
% 組み合わせ
% -------------------------------------------------------------------
\section{装飾の組み合わせ}

\textbf{\textit{太字 + イタリック}}

\textbf{\textcolor{red}{太字 + 赤}}

{\Large \textbf{\textcolor{blue}{大きい + 太字 + 青}}}

\end{document}